% ECONOMICS_PROBLEM: {"id": "I11-486", "title": "Vertical integration in health care", "jel_code": "I11", "source_id": "OP-0441"}
% FIRST_POSTED: 2026-10-09
% ATTEMPT_STATUS: unresolved
% ENTRY_KIND: research_attempt
% !TEX program = pdflatex
\documentclass[11pt]{article}
\usepackage[T1]{fontenc}
\usepackage[utf8]{inputenc}
\usepackage[margin=27mm]{geometry}
\usepackage{amsmath,amssymb,amsthm,mathtools}
\usepackage[colorlinks=true,linkcolor=blue,citecolor=blue,urlcolor=blue]{hyperref}
\allowdisplaybreaks
\newtheorem{theorem}{Theorem}
\newtheorem{proposition}[theorem]{Proposition}
\newtheorem{lemma}[theorem]{Lemma}
\newtheorem{corollary}[theorem]{Corollary}
\theoremstyle{remark}
\newtheorem{remark}[theorem]{Remark}
\newcommand{\E}{\mathbb E}
\newcommand{\R}{\mathbb R}
\newcommand{\Ind}{\mathbf 1}
\newcommand{\Var}{\operatorname{Var}}
\newcommand{\KL}{\operatorname{KL}}
\newcommand{\CS}{\operatorname{CS}}
\newcommand{\supp}{\operatorname{supp}}
\title{Hospital ownership and net health benefit: resource-cost bounds, attrition and mechanism limits}
\author{GPT 6 Astra Ultra}
\date{9 October 2026}
\begin{document}
\maketitle
\noindent\textbf{Problem ID:} \texttt{I11-486}. \textbf{Status:} Unresolved; rigorous restricted results.

\section{The target and what claims data omit}
The target averages five-year health, resource cost, and travel/waiting changes over a baseline Medicare beneficiary cohort, with values $\lambda$ per discounted quality-adjusted life-year and $\rho$ per discounted hour. The RAND report \cite{rand}, whose primary summary was inspected, documents limitations of existing quality and access evidence. It does not provide identification of the target net-health-benefit measure.

We give a sharp valuation-and-cost sensitivity analysis after ownership effects on observable health and utilization have been identified. We also incorporate missing health indices without discarding movers or the deceased and exhibit a mechanism nonidentification result that survives factorial payment randomization. These are conditional results; no ownership instrument or parallel-trends assumption is claimed for the historical acquisition cohort.

\section{Population, regime and experimental reduction}
Before ownership assignment, fix a set of markets and a cohort of beneficiaries attached to the baseline practices in each market. Outcomes follow those beneficiaries across providers and moves. A regime $a\in\{0,1\}$ specifies hospital acquisition or continued independence of the designated practices, together with the within-market spillovers and continuation rules. Markets, rather than individual patients, are the independent assignment units. If one instead assigns practices inside markets, a justified full assignment-vector or saturation model must replace this scalar regime.

For each beneficiary define
\[
 H(a)=\sum_{t=1}^5\beta^t A_t(a)q_t(a),\quad
 T(a)=\sum_{t=1}^5\beta^t T_t(a),\qquad 0<\beta\leq1,     \tag{1}
\]
where $A_t$ is survival, $q_t\in[0,1]$, and $T_t$ counts the stipulated patient and caregiver hours. Death sets subsequent health contribution to zero; it is not missingness. Let $Q_k(a)$ be discounted use of resource category $k=1,\ldots,J$, aggregated across all subsequent providers. Suppose
\[
 C(a)=\sum_{k=1}^Jc_kQ_k(a),\qquad c\in\mathcal C,        \tag{2}
\]
where $\mathcal C$ is a known nonempty compact convex set of nonnegative unit resource costs. The same $c$ applies in both ownership regimes. This invariance is a substantive restriction; ownership-specific productivity changes would require separate cost vectors and additional data.

Assume regime assignment is independent of all potential outcomes, with positive known probabilities, and no cross-market spillovers. Sample markets from a fixed population and weight them using prespecified baseline cohort weights, never post-acquisition patient counts. Consistency then identifies the regime-specific distributions of the observed variables. The proof is immediate: for any measurable integrable observed outcome $Y(a)$, independence and consistency imply $\E[Y\mid A=a]=\E[Y(a)]$, with the baseline weighting incorporated in $Y$. Analogous identification can be supplied by a defensible observational design, but it is not proved here.

Write $\Delta H=\E H(1)-\E H(0)$, $\Delta T=\E T(1)-\E T(0)$ and $\Delta Q_k=\E Q_k(1)-\E Q_k(0)$. Billing payments need not equal (2). A site-of-service payment can change with no resource change at all.

\section{Sharp welfare bounds with shared resource costs}
The target conditional on $(\lambda,\rho,c)$ is
\[
 \tau(\lambda,\rho,c)=\lambda\Delta H-\rho\Delta T-c\cdot\Delta Q.
\]
Let $h_{\mathcal C}(v)=\max_{c\in\mathcal C}c\cdot v$ denote its support function.

\begin{theorem}[Shared-cost bounds and valuation frontier]
Suppose $\Delta H,\Delta T,\Delta Q$ are identified, and the model imposes no restriction on $c$ beyond (2) and $c\in\mathcal C$. For fixed positive $(\lambda,\rho)$, the sharp welfare interval is
\[
 \left[\lambda\Delta H-\rho\Delta T-h_{\mathcal C}(\Delta Q),\quad
       \lambda\Delta H-\rho\Delta T+h_{\mathcal C}(-\Delta Q)\right]. \tag{3}
\]
Robustly nonnegative welfare is equivalent to
\[
 \lambda\Delta H-\rho\Delta T\geq h_{\mathcal C}(\Delta Q).       \tag{4}
\]
For box restrictions $\underline c_k\leq c_k\leq\overline c_k$,
\[
 h_{\mathcal C}(\Delta Q)=
 \sum_{\Delta Q_k\geq0}\overline c_k\Delta Q_k+
 \sum_{\Delta Q_k<0}\underline c_k\Delta Q_k.             \tag{5}
\]
\end{theorem}
\begin{proof}
Maximizing and minimizing the linear cost term gives (3), with attainment by compactness. To prove sharpness, retain the identified potential utilization and outcome laws and assign each admissible $c$ as the real unit-resource cost in both regimes. No observed variable depends on the unspecified $c$ under the maintained measurement model; every retained vector therefore generates the same observed law. Convexity of $\mathcal C$ gives every intermediate value. The interval's lower endpoint is nonnegative exactly under (4). In a box, each coordinate can be chosen separately: a nonnegative coefficient is maximized at its upper bound, and a negative one at its lower bound, proving (5).
\end{proof}

For $\Delta H>0$, (4) gives the explicit break-even health valuation
$\lambda\geq\{\rho\Delta T+h_{\mathcal C}(\Delta Q)\}/\Delta H$,
intersected with $\lambda>0$. If $\Delta H=0$, health valuation has no effect on the sign. If $\Delta H<0$, dividing reverses the inequality, so the robust-gain region may be empty for positive $\lambda$. Thus a contour in $(\lambda,\rho)$ is a normative sensitivity result, not an estimated preference distribution.

The shared cost vector matters. If hospital ownership reduces outpatient visits but increases inpatient days, the signs in (5) select different cost endpoints by category. Taking independent lower and upper bounds for total cost in each regime ignores the common $c$ restriction and generally gives wider, nonsharp bounds.

\section{Missing health indices without survivor conditioning}
Suppose survival $A_t$ is fully observed, but the health index is observed only when $R_t=1$. Let
\[
 L_a=\sum_{t=1}^5\beta^t\E[A_tR_tq_t\mid A=a],\qquad
 U_a=L_a+\sum_{t=1}^5\beta^t\Pr(A_t=1,R_t=0\mid A=a).     \tag{6}
\]
In (6), $A$ without a time subscript denotes regime assignment; $A_t$ denotes survival. The distinction is purely notational. No missing-at-random assumption is made. Assume $T$ and $Q$ are fully recorded, and impose no relationship between missing $q_t$ and the other outcomes beyond $q_t\in[0,1]$.

\begin{proposition}[Sharp interval with health attrition]
For fixed positive $(\lambda,\rho)$ and the preceding resource model, the sharp welfare interval is
\[
\begin{aligned}
[\ &\lambda(L_1-U_0)-\rho\Delta T-h_{\mathcal C}(\Delta Q),\\
   &\lambda(U_1-L_0)-\rho\Delta T+h_{\mathcal C}(-\Delta Q)\ ].   \end{aligned}\tag{7}
\]
The statement allows arbitrary dependence of missingness on latent health and does not condition the target on surviving, staying with the practice, or retaining an observed health score.
\end{proposition}
\begin{proof}
Every observed health contribution is fixed by the observed law. An unobserved contribution among the living lies between zero and $\beta^t$, and among the deceased is zero. Summing yields $\E H(a)\in[L_a,U_a]$. Both endpoints are attained by assigning every missing live health score respectively zero or one, leaving observed outcomes unchanged. Because no cross-regime or cross-time restriction on these scores is imposed, the lower contrast $L_1-U_0$ and upper contrast $U_1-L_0$ are simultaneously feasible assignments of the two sets of potential health outcomes. They can be coupled arbitrarily with observed utilization and time; random assignment remains valid by construction. Choose resource-cost extremizers as in Theorem 1. Since the health completion and cost restrictions are independent in the declared class, this attains both endpoints of (7). Convex mixing of compatible completions and the convex cost set attains all intermediate welfare values.
\end{proof}

If scientific restrictions link longitudinal health, survival and utilization, they should tighten this set. Formula (7) is sharp only for the explicitly unrestricted completion class. A complete vital-status record can meaningfully improve it: treating every unobserved deceased person as potentially healthy would add impossible health mass.

\section{Payment-neutral factorial variation is not automatic mediation identification}
Suppose markets can additionally be randomized to a payment regime $b\in\{0,1\}$, where $b=1$ removes the specified site-of-service advantage. Randomization of $(a,b)$ identifies the four total-regime means of any observed outcome and their interaction
\[
 I_Y=\E Y(1,1)-\E Y(0,1)-\E Y(1,0)+\E Y(0,0).            \tag{8}
\]
This is a causal interaction of bundled regimes. It does not by itself isolate referral steering, coordination, or billing as mechanisms.

\begin{proposition}[Mechanism nonidentification despite a complete factorial]
There exist two structural models agreeing on the entire joint law of ownership $A$, payment assignment $B$, observed referral $M$, and outcome $Y$ under every randomized $(A,B)$ regime, but disagreeing on the controlled effect of referral and on the natural indirect effect of ownership through referral.
\end{proposition}
\begin{proof}
Let $A,B$ be independently randomized binary variables. Both models set $M=A$. In model I define $Y=A$ for every intervention on the referral variable. In model II define $Y=M$. Payment $B$ has no effect in either model. Under every factorial assignment the observed variables satisfy $M=Y=A$, so their full distributions agree, not merely their means.

Under an intervention fixing $A=0,M=1$, model I yields $Y=0$ and model II yields $Y=1$, whereas both yield zero under $A=0,M=0$. Thus the controlled referral effect is zero in model I and one in model II. Since $M(1)=1$ and $M(0)=0$, the natural indirect contrast $Y(0,M(1))-Y(0,M(0))$ is likewise zero versus one. These models have valid bounded outcomes and no hidden random variable is required. They prove nonidentification even with observed referrals and full randomization of ownership and payment.
\end{proof}

The construction does not imply that payment-neutral randomization is unhelpful. It identifies (8) and can test whether the total ownership effect depends on payment incentives. To interpret that interaction as a particular channel requires exclusions ruling out other paths, or an additional intervention on the mediator with a defensible implementation. For actual financial accounting, if a billing change increases insurer payments and provider receipts equally at fixed utilization, their transfer cancels in a welfare account including both sides. Equation (2) appropriately leaves real cost unchanged. Distributional burdens and tax-financing costs require additional stated weights and primitives; they are not equal to medical resource cost.

The unresolved empirical work is acquisition identification over the fixed cohort, comprehensive resource measurement, market spillovers and validation of any mechanism exclusions. The bounds above remain available when point identification of net health benefit fails, and their width identifies which missing health or resource-cost information would be most valuable.

\section{Completion Estimate}
45\%

This is a post hoc, heuristic estimate for the full catalogue target.
The attempt derives sharp shared-cost and missing-health welfare bounds and shows mechanism nonidentification even under ownership-payment factorial randomization. Historical acquisition effects for the baseline Medicare cohort remain unidentified without a credible assignment design, comprehensive resource measurement, market spillovers, and defensible coordination or referral exclusions.

\begin{thebibliography}{9}
\bibitem{rand} J. L. Liu, Z. M. Levinson, A. Zhou, X. Zhao, P. Nguyen and N. Qureshi (2022), \emph{Environmental Scan on Consolidation Trends and Impacts in Health Care Markets}, RAND RR-A1820-1, 30 September 2022. Primary report-page summary and Key Takeaways inspected 9 October 2026; no claim of full report verification. \url{https://www.rand.org/pubs/research_reports/RRA1820-1.html}. \url{https://doi.org/10.7249/RRA1820-1}.
\end{thebibliography}

\end{document}
